%% hopfion-framework/electroweak/yukawa.tex
%% Sources: main_paper4.tex, main_paper1.tex
%% Yukawa couplings, MSbar-to-pole scheme conversion, tracking relation
%% Auto-generated from project files. Edit source papers to update.
%%
%% Status tags: \status{published}, \status{open}, \status{conjectured}
%% \source{PaperN, label} — where the proof lives
%% \residual{X\%} or \residual{exact} — numerical accuracy


%════════════════════════════════════════════════════════════════════
% Verlinde Cancellation — Paper IV
%════════════════════════════════════════════════════════════════════
\begin{theorem}[Verlinde cancellation~\cite{Verlinde1988,Paper3}]
\label{P4:thm:verlinde_cancel}
For all $R_0 > 0$, the IR virial fixed point satisfies
\begin{equation}
  V^*(R_0) = \vp \quad \text{exactly.}
  \label{P4:eq:Vstar_exact}
\end{equation}
The $R_0$-dependent correction factors in $V(0;\,R_0)$
and in $V^*/V(0;\,R_0)$ cancel identically.
\end{theorem}
\status{published}
\source{P4:thm:verlinde_cancel}
\residual{exact}

\begin{corollary}[Verlinde spoke amplitude is exact for all $R_0$]
\label{P4:cor:verlinde_exact}
The factor $\vp^{2Q} = \vp^{20}$ in the electron mass formula
requires no finite-$R_0$ correction.
The deviation $V^*(R_0{=}3) = 1.61824 \neq \vp = 1.61803\ldots$
observed in the numerical solver at grid spacing $h = 0.12$
is a discretisation artefact of order $O(h^2)$, not a physical
finite-torus effect.
At $h=0.12$: the per-sector numerical noise is $0.013\%$,
accumulating to $0.26\%$ over $2Q=20$ Verlinde sectors ---
consistent with $O(h^2) = 0.0144$ scaling.
\end{corollary}
\status{published}
\source{P4:cor:verlinde_exact}
\residual{exact}

%════════════════════════════════════════════════════════════════════
% WZW Amplitude, Koide, and Yukawa — Paper I
%════════════════════════════════════════════════════════════════════
\begin{proposition}[Eta identity]
    \label{P1:prop:eta}
    \begin{equation}
      \left[\frac{\eta(\vp\,i)}{\eta(i/\vp)}\right]^{\!2}
      \;=\; \frac{1}{\vp}
      \;=\; \frac{S_{0,0}}{S_{0,1/2}}\bigg|_{k=3}.
      \label{P1:eq:eta_identity}
    \end{equation}
    \end{proposition}
\status{published}
\source{P1:prop:eta}

\begin{proposition}[Cardy boundary ratio~\cite{Cardy1989}]
    \label{P1:prop:cardy}
    In the $\mathrm{SU}(2)_3$ WZW model, the cylinder partition
    functions with Cardy boundary representations $j_A$, $j_B$ are
    $Z_{j_A,j_B}(\tau)=\sum_l(S_{j_A,l}S_{j_B,l}/S_{0,l})\chi_l(\tau)$.
    As $\mathrm{Im}(\tau)\to\infty$:
    \begin{equation}
      \frac{Z_{0,0}(\tau)}{Z_{0,1/2}(\tau)}
      \;\longrightarrow\;
      \frac{S_{0,0}}{S_{0,1/2}}
      \;=\; \frac{1}{\vp}.
      \label{P1:eq:cardy_ratio}
    \end{equation}
    \end{proposition}
\status{published}
\source{P1:prop:cardy}

\begin{theorem}[Koide $K=2/3$ automatic]
\label{P1:thm:koide}
For any three masses parametrised by the $\mathbb{Z}_3$-symmetric ansatz
\begin{equation}
  \sqrt{m_g} = M\!\left(1 + \sqrt{2}\cos\!\left(\theta+\frac{2\pi(g-1)}{3}\right)\right),
  \label{P1:eq:koide_ansatz}
\end{equation}
the Koide ratio $K = \sum m_g/(\sum\sqrt{m_g})^2 = 2/3$ exactly,
for all $\theta$ and $M$.
\end{theorem}
\status{published}
\source{P1:thm:koide}
\begin{proof}
$\mathbb{Z}_3$ orthogonality gives $\sum m_g = M^2(3 + 2\sum\cos(\ldots)) = 6M^2$
and $(\sum\sqrt{m_g})^2 = 9M^2$, so $K=6/9=2/3$.
\end{proof}
\status{published}
\source{P1:thm:koide}
\residual{exact}

\begin{result}[Electron Yukawa]
\label{P1:res:yukawa}
The electron Yukawa coupling is the WZW amplitude on the Hopfion torus
with charge $Q=10$:
\begin{equation}
  \boxed{y_e \;=\; \exp\!\bigl(-2\pi Q\,T_1\bigr)
  = \exp\!\left(-\frac{25\pi}{6}\right) \;\approx\; 2.066\times10^{-6}.}
  \label{P1:eq:yukawa}
\end{equation}
This gives $v_{\mathrm{EW}} = m_e/y_e \approx 247.4\;\text{GeV}$
(deviation $0.46\%$ from $246.2\;\text{GeV}$).
The internal consistency check $25\pi/6\times 12/(5\pi) = 10 = Q$ holds exactly.
\end{result}
\status{proved (Paper~II)}
\source{P1:res:yukawa}
\residual{0.46\%}
\depends{foundations/wzw_group.tex (P1:thm:charge_q10, P1:thm:k3)}

\begin{result}[Koide angle from WZW holonomies]
\label{P1:res:koide_angle}
The Chern--Simons phases for generation $g$ are
$\phi_g^{\mathrm{CS}} = 2\pi k_g\cdot\tfrac{3}{4}/(k_g+2)$,
giving $\phi_1^{\mathrm{CS}} = \pi/2$, $\phi_2^{\mathrm{CS}} = 3\pi/4$,
$\phi_3^{\mathrm{CS}} = 9\pi/10$.
The predicted Koide angle is
\begin{equation}
  \theta_K \;=\;
  \arg\!\left(
    1\cdot e^{i\pi/2}
    + 2^{1/4}\cdot e^{i3\pi/4}
    + \vp^{1/2}\cdot e^{i9\pi/10}
  \right) = 132.550^{\circ}
  \label{P1:eq:thetaK}
\end{equation}
vs.\ observed $\theta_K^{\mathrm{exp}} = 132.732^{\circ}$ (deviation $0.14\%$).
\end{result}
\status{proved (Paper~II)}
\source{P1:res:koide_angle}
\residual{0.14\%}

